\chapter[Sermon 43. Christ a Strong Angel, Is Set Against Antichrist: And Is Excellently Described, Shining Again in the Darkness of the Church with Consolation.]{Christ a Strong Angel, Is Set Against Antichrist: And Is Excellently Described, Shining Again in the Darkness of the Church with Consolation.}
\label{sermon:043}
\markright{Sermon 43. Christ a Strong Angel, Is Set Against Antichrist: And ...}

And I saw another mighty angel come down from Heaven, clothed with a cloud, and the rainbow upon his head, and his face was as the sun, and his feet as it were pillars of fire. And he had in his hand a little book open. And he put his right foot upon the sea, and his left foot on the earth. And he cried with a loud voice, as when a lion roars. And when he had cried, seven thunders spoke their voices. And when the seven thunders had spoken their voices, I was about to write. And I heard a voice from Heaven saying unto me: Seal up those things which the seven thunders spoke, and write them not.

\index{Antichrist}\index{Mahomet}\index{Turks}\index{Papacy}Hitherto we have heard many things about the most dangerous conflicts against Christian piety and truth. But nothing has been spoken of the defenders and maintainers of the true religion; rather, the success and wonderful felicity of the wicked has been preached, especially in the fifteenth and sixteenth trumpets, namely, under Papistry and Mahometism. Therefore, it might seem to many that truth itself was not only oppressed and lost, but also that the truth of God’s promises began to faint and decay. For the godly are oppressed; iniquity triumphs in all places; truth is trampled underfoot, and lawlessness reigns everywhere. Who, therefore, would not think that the things spoken everywhere about the rewards of the good and the punishment of the evil are in vain? Doubtless, the children of Israel doubted the faith of God’s promises when the term of their captivity was extended only to seventy years. What marvel is it, then, if the faithful at this day also, seeing the servitude or bondage of Mahomet and the tyranny of the Pope, or Antichrist, has continued for many years, begin also to be tempted, as also the saints were tempted in old times, witness Asaph, Psalm 73? And truly, there will be those at this day who will say, “Perhaps this world will always be, will never have an end; Papistry will reign forever; the Mahometans will conquer forever; the godly will be miserable forever; therefore it is better to conform to the world.” Our fathers, eight hundred and a thousand years ago, have long looked for your judgment, and thought the last day of the Lord was at hand, but no end appears anywhere, but all things are daily renewed; therefore, the same face of the world will always be, the courses of time duly returning. Who, therefore, knows whether a reward is prepared for the godly, or punishment for the wicked? For far other things chance to holy men than many look for. They look for blessing, life, and felicity, and behold they are overwhelmed with the curses of all men, carried to execution, and seem the most unfortunate of all men. He who cleaves to the Papists, Turks, and enemies of the Gospel, goes through luckily enough, and so forth.

Therefore, as everywhere in this book, joyful things are mixed with sorrowful, so here also, after the most grievous battles of the fifth and sixth trumpets and the most strong temptations, he joins a most joyous gospel for the consolation and confirmation of the faithful, lest they should anywhere doubt of God’s promises or revolt from the true religion to the false. Therefore, against Antichrist, the black angel of the bottomless pit, is set the bright or shining angel of heaven, the Lord Christ. Here is gallantly described and is said to return unto his; the same swears solemnly that there shall be no other time, but that in the seventh trumpet the very mystery of God should be fulfilled.

Moreover, the Lord Christ commands John to eat the open book which Christ held in his hand and to prophesy again. By all these things, for the comfort of all the godly, is signified that Christ shall return into the church, out of which he seemed by his enemy and vicar to be cast out, with great glory and power; neither the hope and expectation of the faithful to be vain, however the last day of judgment be deferred into many ages, and the godly feel great adversity. Finally, that punishment and reward is prepared by God, and that this shall be given to the godly, and that inflicted to the wicked. For to the intent we might hereof be most assured, Christ takes a solemn oath and says it shall come to pass that the catholic and Christian truth shall again come into the field and fight valiantly against the Antichristian and Mahometan doctrine. Hereof, therefore, shall we learn not to despair in the long persecutions of Antichrist and Mahomet. We shall learn also how to fight against Antichrist and how he must be overcome not with warlike but spiritual weapons. These he is not able to match. He whets one sword on another. And hitherto, in deed, in these two last chapters, has been spoken of the wars of heretics and of the ungodly, and of Antichrist, the head of all evil; hereafter shall follow of the contrary fight of the godly and the maintenance of godliness.

\index{John, Saint}Before this is set a description of Christ most elegant, most wholesome, and most full of consolation, declaring his power in the ministry, by the ministers of the word, whom he has clothed with authority from above, and by the weak things of this world, overcomes and defeats the strongest things of this world, which seemed invincible. Christ, king and bishop, animates all his faithful with his spirit and word equally, always and everywhere works many things by his power, so that he is now felt by all to have come again, unto whom he seemed hitherto to have absent himself somewhat long. And I doubt whether in all the canonical books, in the prophecy of Isaiah, after the story of the Gospel, and especially after the Gospel of blessed Saint John, there be any other book which has more and more goodly descriptions of Christ than this book. They are deceived and much abused who suppose a rare gospel to be preached in this book. But let us see the description of Christ by parts.

\index{Hebrews, Epistle to the}We have shown in the fifth and sixth verses that Antichrist, the Pope and Mahomet, are powerful. Now, a mighty angel, the Lord Christ himself, stands against them; an angel in deed, but not in nature or dignity. For he took the seed of Abraham, not the nature of an angel, and is much greater than angels, though he is Lord of angels. Which the Apostle shows in the first and second chapters to the Hebrews. But it is the angel of great counsel, namely, the ambassador of God the Father unto us, as Isaiah and Malachi called him, to teach the will of the heavenly Father, and to work our salvation, and now also appointed that from the right hand of the Father, as king and bishop, he should keep and defend his church. This Christ, I say, given to men, is strong, not weak. Strong to overcome and break through all the force of this world, of his enemies, of Antichrist, of the Devil, and of hell gates. No man therefore need to doubt that he may be defended by this strong giant; no man need to despair in any perils or matters, be they ever so desperate, since Christ almighty lives, who is able to help weary matters. He therefore must be called upon by us in all dangers; of him we must look for help patiently and with steadfast faith.

A star has indeed fallen from heaven, but in the time of Antichrist, something has crept out of the bottomless pit, and Christ descends from heaven. The Lord will not come from heaven in a physical body until the final judgment, but he is said to return spiritually whenever he appears to have withdrawn or absented himself. For he is never absent; he is always present with his people, who are everywhere. He is said to be absent when his help is not felt, but people are troubled by adversity and seemingly broken by evil. Christ therefore descends to us whenever he gives us aid and counsel when we are in need. So, in the most dangerous conflicts of Antichrist and Mahomet, and finally of heretics and all other adversaries, Christ himself with his spirit is present with his people, helps them, encourages and strengthens them, and defends them.

Antichrist, as it were clothed and wrapped with the smoke of the bottomless pit, is evil favored. And Christ covered and as it were decked with a cloud, both shows himself to be God, who can convey him above the air and clouds into his kingdom, and can rain down dew or wholesome showers to water and make fruitful his chosen. A cloud many times in the holy history was a token of God present. A cloud took up Christ as it were girded from the eyes of his disciples. Clouds with Isaiah in the 45th and 60th chapters drop down grace. Finally, we shall be taken up in the clouds with all the chosen to meet the Lord in the air. Antichrist is crowned with a crown, which he obtained for himself by crafty means, by flattery, threatenings, and deceits, and keeps it by wicked practices, bloody weapons, and all kinds of crafty dealings and ungodliness. But the rainbow is on the head of Christ. For he is king of peace, pacifying and reconciling all things unto God, reconciling also themselves together with the bond of charity. The rainbow is a token of the league and amity of God towards us, as we mentioned before in the matters of Noah. Of Antichrist is woe, desperation, anguish, and a most afflicted and troubled conscience. Christ is a consolation and peace of the conscience, so that no man need to say, “The mountains fall upon us and cover us.” For the faithful, being delivered, cry, “Abba, Father.” The same Lord Christ said also in the Gospel, “In the world you have affliction, but in me peace.” And again, “Come to me all that labor and are burdened, and I will refresh you.” Therefore, the rainbow sits on no man’s head more rightly than upon Christ’s head. For in the head of this angel appeared the rainbow, as though it had been his crown or diadem. And now we understand from whence consolation and pacification is to be looked for in most grievous perils and afflictions of Antichrist.

The sun, seen through the smoke of the bottomless pit, raised by the angel thereof and darkened most filthily, was a great grief to the world. But the face of this our angel is clear, and clarifies, and shines as bright as the sun, as in times past also he is said to have shone in Matthew 17. The bright face of Christ brings joy and unspeakable gladness to those who behold it, and pacifies the minds. And this is seen by us spiritually and by faith. Paul is the best expounder of this passage in 2 Corinthians 3 and 4. Certainly, full knowledge of Christ, whereby we believe him to be given to us by the Father, that by him all things should be accomplished, and he himself alone should be everything to us, arises as a most joyful light in the minds. For we are enlightened by his light, as is declared in the gospel of John. However, the terrible darkness of errors and calamities may be overcast in the world by the prince of darkness through the pope and Mahomet, and other corrupters of God’s doctrine, yet the light which is in the minds of the faithful through faith in Christ directs, lights, comforts, and preserves.

Antichrist has the tails of scorpions, most venomous; but the feet of the Lord Christ are as pillars of fire. All things of Christ are firm, right, and most clean. In another place also, the truth of the Lord is figured by pillars. And God himself is called a consuming fire, burning up all uncleanness. Moreover, fire pierces, and cannot be easily quenched when it begins to catch and to burn. And who shall restrain the course of the fire of the Holy Spirit? With this is joined another thing, that the Angel set his right foot upon the sea, and his left on the earth. But to set foot is to challenge possession to himself. For as we say in Dutch, to kindle the mind to dwell in any place; so we read of the Hebrew phrase, whatever your foot shall tread upon, shall be yours: that is, whatever he shall win, shall be your own right and possession. And Christ seemed as it were to be cast out of his possession by Antichrist, Mahometans, and the rest of his enemies. For we have heard many times that they have won lands and islands. But this vision instructs us that Christ recovers again, reenters possession, and brings under his subjection such places which seemed to have been taken and lost both by sea and land, throughout all parts of the world. For by the preaching of the gospel, many are recovered who were plucked away from the true religion; that now I need not rehearse that Christ has in all parts of the world his chosen, who never bowed their knees to Baal. In them he has a most strong and most purified kingdom, figured by the fiery pillars. Certainly, the Apostle calls the church a pillar and foundation of truth. 1 Timothy 2. And Christ himself also says that his kingdom in the world is immovable, however it be assailed by Antichrist. The gates of hell shall not prevail against it, and his church.

And just as Antichrist desires to have the book of the gospel shut, and closed, obscure and imperfect, so the Lord Christ holds the book open in his hand. He opens, and no man shuts. You understand therefore from what it arises that Antichrist, although he be of most great power, cannot at this day shut the gospel book, which he seeks with all his force to do. Of Christ and his Spirit we have the word bright and clear. By the grace of Christ we have the bright preaching of the gospel, contrariwise a dark and an intricate sophistry of Antichrist: of this book shall follow hereafter more plentiful things. Hereunto appertains the worthy invention and godly benefit of Printing, never commended enough. This opens books, and sends them abroad into the world in spite of all the enemies of God’s truth, and scatters them abroad in every corner of the world. So that they which cannot hear preachers, to them come godly books not without fruit.

\index{Erasmus of Rotterdam}The sound of Antichrist’s wings is like the noise of chariots when many horses run into battle; therefore Christ also makes a noise, and cries indeed with a loud voice. John adds a parable. He roars—or loweth—which is as much as to say, he roars. For Erasmus, perhaps, says that the Greeks do not keep the difference between \textit{brysō} and \textit{goggysō} as the Latins do between \textit{rugire} and \textit{mugire}, that is, to roar and to low. He roars as a lion. We have heard before how Christ is called a lion of the tribe of Judah. Therefore, just as when a lion rears, as Amos notes in chapter 3, all are afraid of themselves; so when the Lord Christ cries by his word, all the wicked are assembled. This signifies that the gospel shall be preached against Antichrist constantly and with authority, to the terror of God’s enemies. And without doubt, although the princes, both spiritual and temporal, seem to despise and utterly reject the preaching of the gospel, it is certain by many tokens and conjectures that they are exceedingly afraid of that preaching, which they despise as vile. For they endeavor with all their strength of mind and apply all their counsels to abolish that same preaching. But if they suppose it to be of no force, why are they so afraid of it? Why are they at such great expense? Why can they never be at rest? To this pertains the common prophecy, which tells that while Antichrist reigns, Elijah shall come, who with a sharp, lively, and manly preaching shall refute the triflings of Antichrist. Therefore, the spirit of Elijah and his earnest preaching is that roaring of the lion, which roars out Christ’s truth.

And as soon as that roaring was heard, the thunders also uttered their voices. By which voices are signified the sundry graces of the Holy Spirit, and chiefly the terrible preaching of the truth of the canonical scriptures, as appeared in chapter 4 of this book. For with the gospel in the latter ages, the scripture of the prophets shall be expounded again, which seems as it were to thunder against Antichrist, against sins, and wicked people. Verily James and John, brethren and Apostles of Christ, are called by the Lord Boanerges, that is, the sons of thunder, thunderers, that is to wit, excellently sharp in preaching, and to be feared.

\index{Daniel (Book of)}And St. John would by and by have written the voices of the thunders, but he is forbidden so to do; he is commanded only to seal them. For since the holy scripture, through the inspiration of the Holy Spirit, was written and brought forth already by the servants of God, the Prophets and Apostles, what need were there to write and set it forth again? Those things are sufficient for the godly that are written. From these, preachers may take thundering, lightning, and thunderbolts against Antichrist and all sects. And where he is commanded to seal up the things already set forth, it alludes to the last chapter of Daniel; and that sealing is referred as well to the godly as to the ungodly. Sealed letters are doubtless of most authority. St. John therefore, by his sealing, does not now make the Scriptures authentic, but in sealing them declares that they are authentic enough; so to the scriptures no godly man goes about either to add or diminish. Thus I say, the scriptures are sealed to the godly, as to those who are most persuaded that the scriptures are most perfect and authentic, which may most fully suffice in the church for true piety against all ungodliness. Where the wicked will not see this, and seek not for all things of godly religion in the scriptures, neither care greatly for the scriptures: what marvel is it, though the scriptures be sealed to them, that is shut up, which they neither greatly care for, nor understand, nor yet will understand? And on this wise is Christ set against Antichrist, and recovers again his church, discomfits and subdues Antichrist: to whom be malediction forever.
